\begin{lemma}
Let \(G\) be a finite \(\Delta\)-regular graph, let \(0<\gamma\le \Delta\),
let \(\nu\) be a law on activation--priority pairs \((A,\pi)\), and let
\(\mu\) be a nonnegative mass function on subsets of \(V(G)\).  Assume that
\(\nu\) has total mass one; that every
activation atom has mass
\[
\nu(A=A_0)=
\left(\frac{\gamma}{\Delta}\right)^{|A_0|}
\left(1-\frac{\gamma}{\Delta}\right)^{|V(G)|-|A_0|};
\]
that priorities on each activation atom are supported on the unit cube; that
for every \(t\in[0,1]^{V(G)}\),
\[
\nu(A=A_0,\ 0\le \pi(z)\le t_z\text{ for all }z)
=\nu(A=A_0)\prod_{z\in V(G)}t_z;
\]
that, for every \(z\in A_0\), the fixed-activation comparison event
\[
A=A_0,\quad \pi(q)\in[0,1]\text{ for all }q\in A_0,\quad
\pi(q)<\pi(z)\text{ for all }q\in A_0\cap N_G(z)
\]
has mass
\[
\nu(A=A_0)\int_0^1 a^{|A_0\cap N_G(z)|}\,da;
\]
and that \(\mu\) is the push-forward of \(\nu\) by the output map
\[
I(A,\pi)=\{z\in A:\pi(z)>\pi(w)\text{ for every }w\in A\cap N_G(z)\}.
\]
If \(u\ne v\) and \(u,v\) are nonadjacent, then
\[
\sum_S {\bf 1}_{\{u,v\}\subseteq S}\mu(S)
=
\frac{2}{\Delta^2}\int_0^\gamma\int_x^\gamma
\left(1-\frac{x}{\Delta}\right)^{
\Delta-|N_G(u)\cap N_G(v)|}
\left(1-\frac{y}{\Delta}\right)^\Delta\,dy\,dx .
\]
\end{lemma}
\begin{proof}
Let
\[
E_S=\{(A,\pi): I(A,\pi)=S\}.
\]
The events \(E_S\) are pairwise disjoint as \(S\) ranges over all subsets of
\(V(G)\).  By the push-forward hypothesis,
\[
\operatorname{ofReal}(\mu(S))=\nu(E_S)
\]
for every \(S\).  Since \(\mu(S)\ge 0\), this identity is an identity of the
underlying real masses: no negative value can be hidden by the
\(\operatorname{ofReal}\) map.  Hence, by finite additivity over the disjoint
subfamily indexed by the sets \(S\) containing both \(u\) and \(v\),
\[
\sum_S {\bf 1}_{\{u,v\}\subseteq S}\mu(S)
=\nu\{(A,\pi):u\in I(A,\pi)\text{ and }v\in I(A,\pi)\}.
\]
This is the only place where the nonnegativity of \(\mu\) is used.

It remains to compute this last probability from the activation and priority
laws.  This computation is exactly
\noderef{SamplingNonadjacentPairSurvivalChamberLaw}, applied with the present
graph \(G\), parameters \(\Delta,\gamma\), law \(\nu\), and vertices \(u,v\).
Its hypotheses are precisely the activation atom formula, unit-cube support,
lower-rectangle product law, fixed-activation comparison identity,
\(\Delta\)-regularity, \(0<\gamma\le\Delta\), and the nonadjacency of \(u\)
and \(v\), all of which are among the present assumptions.  Therefore
\[
\nu\{(A,\pi):u\in I(A,\pi)\text{ and }v\in I(A,\pi)\}
=
\operatorname{ofReal}\!\left[
\frac{2}{\Delta^2}\int_0^\gamma\int_x^\gamma
\left(1-\frac{x}{\Delta}\right)^{
\Delta-|N_G(u)\cap N_G(v)|}
\left(1-\frac{y}{\Delta}\right)^\Delta\,dy\,dx\right].
\]
The range hypotheses \(0<\gamma\le\Delta\) imply that
\(0\le x/\Delta\le1\) and \(0\le y/\Delta\le1\) throughout the integration
chamber \(0\le x\le y\le\gamma\).  Hence the integrand and the prefactor are
nonnegative, so this extended-nonnegative identity is the same real mass
identity used in the preceding finite-additivity calculation.
Combining this identity with the finite push-forward identity proved above
gives the displayed formula for
\(\sum_S {\bf 1}_{\{u,v\}\subseteq S}\mu(S)\).
\end{proof}
